%LTeX: language=fr-FR
\documentclass[a4paper,11pt]{article}

% ---------- Préambule minimal ----------
\usepackage[french]{babel}
\usepackage{amsmath,amssymb}
\usepackage{enumitem}
\usepackage{tikz}
\usepackage{tasks}
\usetikzlibrary{arrows.meta,patterns}
\usepackage[margin=2cm]{geometry}
\usepackage{needspace}
\usepackage{currfile}
\usepackage{fancyhdr}
\usepackage{lastpage}
\usepackage{array}
\settasks{label=\alph*.}
\setlength{\parindent}{0pt}

\pagestyle{empty}
\input{\currfiledir ../def.tex}
\def \Document {Exercices}


% Environnement "exercice" : \begin{exercice}{<numéro>} ... \end{exercice}
\newenvironment{exercice}[1]{%
  \Needspace{6\baselineskip}\par\bigskip\noindent\textbf{Exercice #1}\par\smallskip\noindent\ignorespaces
}{%
  \par
}

\newcolumntype{C}[1]{>{\centering\arraybackslash}m{#1}}

% ---------- Repères et courbes ----------
% Quadrillage + axes : \axes{xmin}{xmax}{ymin}{ymax}
\newcommand{\axes}[4]{%
  \draw[gray!35, very thin, step=1] (#1,#3) grid (#2,#4);
  \draw[-{Latex[length=2mm]}] (#1,0) -- (#2+0.4,0);
  \draw[-{Latex[length=2mm]}] (0,#3) -- (0,#4+0.4);
}

% Exercices 1 et 6 : fonction f définie sur [-7 ; 7]
\newcommand{\repereUn}[1]{%
  \begin{tikzpicture}[x=0.55cm, y=0.55cm]
    \axes{-7}{7}{-5}{5}
    \foreach \x in {-7,...,-1,1,2,...,7} \node[below] at (\x,0) {\scriptsize $\x$};
    \foreach \y in {-5,...,-1,1,2,...,5} \node[left] at (0,\y) {\scriptsize $\y$};
    \node[below left] at (0,0) {\scriptsize $0$};
    \draw[very thick] plot coordinates {(-7.00,-2.000) (-6.90,-1.809) (-6.80,-1.632) (-6.70,-1.463) (-6.60,-1.296) (-6.50,-1.125) (-6.40,-0.944) (-6.30,-0.747) (-6.20,-0.528) (-6.10,-0.281) (-6.00,0.000) (-5.90,0.307) (-5.80,0.626) (-5.70,0.950) (-5.60,1.277) (-5.50,1.600) (-5.40,1.915) (-5.30,2.218) (-5.20,2.502) (-5.10,2.765) (-5.00,3.000) (-4.90,3.213) (-4.80,3.412) (-4.70,3.598) (-4.60,3.771) (-4.50,3.931) (-4.40,4.079) (-4.30,4.214) (-4.20,4.338) (-4.10,4.449) (-4.00,4.550) (-3.90,4.640) (-3.80,4.718) (-3.70,4.787) (-3.60,4.845) (-3.50,4.894) (-3.40,4.933) (-3.30,4.963) (-3.20,4.984) (-3.10,4.996) (-3.00,5.000) (-2.90,4.957) (-2.80,4.840) (-2.70,4.662) (-2.60,4.440) (-2.50,4.188) (-2.40,3.920) (-2.30,3.653) (-2.20,3.400) (-2.10,3.178) (-2.00,3.000) (-1.90,2.850) (-1.80,2.704) (-1.70,2.564) (-1.60,2.432) (-1.50,2.312) (-1.40,2.208) (-1.30,2.122) (-1.20,2.056) (-1.10,2.014) (-1.00,2.000) (-0.90,2.003) (-0.80,2.013) (-0.70,2.028) (-0.60,2.049) (-0.50,2.074) (-0.40,2.104) (-0.30,2.138) (-0.20,2.175) (-0.10,2.216) (0.00,2.259) (0.10,2.305) (0.20,2.352) (0.30,2.401) (0.40,2.450) (0.50,2.500) (0.60,2.550) (0.70,2.599) (0.80,2.648) (0.90,2.695) (1.00,2.741) (1.10,2.784) (1.20,2.825) (1.30,2.862) (1.40,2.896) (1.50,2.926) (1.60,2.951) (1.70,2.972) (1.80,2.987) (1.90,2.997) (2.00,3.000) (2.10,2.988) (2.20,2.952) (2.30,2.894) (2.40,2.816) (2.50,2.719) (2.60,2.604) (2.70,2.473) (2.80,2.328) (2.90,2.170) (3.00,2.000) (3.10,1.820) (3.20,1.632) (3.30,1.437) (3.40,1.236) (3.50,1.031) (3.60,0.824) (3.70,0.616) (3.80,0.408) (3.90,0.202) (4.00,0.000) (4.10,-0.202) (4.20,-0.408) (4.30,-0.616) (4.40,-0.824) (4.50,-1.031) (4.60,-1.236) (4.70,-1.437) (4.80,-1.632) (4.90,-1.820) (5.00,-2.000) (5.10,-2.170) (5.20,-2.328) (5.30,-2.473) (5.40,-2.604) (5.50,-2.719) (5.60,-2.816) (5.70,-2.894) (5.80,-2.952) (5.90,-2.988) (6.00,-3.000) (6.10,-2.962) (6.20,-2.856) (6.30,-2.694) (6.40,-2.488) (6.50,-2.250) (6.60,-1.992) (6.70,-1.726) (6.80,-1.464) (6.90,-1.218) (7.00,-1.000)};
    \node at (-6.3,3) {$\mathcal{C}_f$};
    #1
  \end{tikzpicture}%
}

% Exercice 3 : repère vide de -4 à 5
\newcommand{\repereVide}[1]{%
  \begin{tikzpicture}[x=0.6cm, y=0.6cm]
    \axes{-4}{5}{-4}{5}
    \foreach \x in {-4,...,-1,1,2,...,5} \node[below] at (\x,0) {\scriptsize $\x$};
    \foreach \y in {-4,...,-1,1,2,...,5} \node[left] at (0,\y) {\scriptsize $\y$};
    \node[below left] at (0,0) {\scriptsize $0$};
    #1
  \end{tikzpicture}%
}

% Exercice 5 : solde du compte (unité en ordonnée : 100 €)
\newcommand{\repereBanque}[1]{%
  \begin{tikzpicture}[x=1cm, y=0.5cm]
    \draw[gray!35, very thin, step=1] (0,-2) grid (12,6);
    \draw[-{Latex[length=2mm]}] (0,0) -- (12.5,0) node[above left] {\scriptsize mois};
    \draw[-{Latex[length=2mm]}] (0,-2) -- (0,6.4) node[above] {solde (€)};
    \foreach \x in {1,...,12} \node[below] at (\x,0) {\scriptsize $\x$};
    \foreach \y/\l in {-2/-200,-1/-100,1/100,2/200,3/300,4/400,5/500,6/600} \node[left] at (0,\y) {\scriptsize $\l$};
    \node[below left] at (0,0) {\scriptsize $0$};
    \draw[very thick] plot coordinates {(0.00,1.000) (0.10,1.150) (0.20,1.299) (0.30,1.447) (0.40,1.592) (0.50,1.734) (0.60,1.873) (0.70,2.007) (0.80,2.136) (0.90,2.259) (1.00,2.375) (1.10,2.484) (1.20,2.584) (1.30,2.675) (1.40,2.757) (1.50,2.828) (1.60,2.888) (1.70,2.936) (1.80,2.971) (1.90,2.993) (2.00,3.000) (2.10,2.958) (2.20,2.840) (2.30,2.663) (2.40,2.440) (2.50,2.188) (2.60,1.920) (2.70,1.653) (2.80,1.400) (2.90,1.177) (3.00,1.000) (3.10,0.860) (3.20,0.736) (3.30,0.627) (3.40,0.528) (3.50,0.438) (3.60,0.352) (3.70,0.269) (3.80,0.184) (3.90,0.095) (4.00,0.000) (4.10,-0.109) (4.20,-0.232) (4.30,-0.363) (4.40,-0.496) (4.50,-0.625) (4.60,-0.744) (4.70,-0.847) (4.80,-0.928) (4.90,-0.981) (5.00,-1.000) (5.10,-0.981) (5.20,-0.928) (5.30,-0.847) (5.40,-0.744) (5.50,-0.625) (5.60,-0.496) (5.70,-0.363) (5.80,-0.232) (5.90,-0.109) (6.00,0.000) (6.10,0.100) (6.20,0.200) (6.30,0.300) (6.40,0.400) (6.50,0.500) (6.60,0.600) (6.70,0.700) (6.80,0.800) (6.90,0.900) (7.00,1.000) (7.10,1.098) (7.20,1.194) (7.30,1.287) (7.40,1.381) (7.50,1.475) (7.60,1.571) (7.70,1.671) (7.80,1.774) (7.90,1.884) (8.00,2.000) (8.10,2.130) (8.20,2.278) (8.30,2.442) (8.40,2.619) (8.50,2.806) (8.60,3.001) (8.70,3.200) (8.80,3.402) (8.90,3.602) (9.00,3.800) (9.10,3.992) (9.20,4.174) (9.30,4.346) (9.40,4.503) (9.50,4.644) (9.60,4.765) (9.70,4.864) (9.80,4.938) (9.90,4.984) (10.00,5.000) (10.10,4.983) (10.20,4.933) (10.30,4.853) (10.40,4.744) (10.50,4.609) (10.60,4.451) (10.70,4.271) (10.80,4.072) (10.90,3.856) (11.00,3.625) (11.10,3.382) (11.20,3.128) (11.30,2.866) (11.40,2.599) (11.50,2.328) (11.60,2.056) (11.70,1.785) (11.80,1.517) (11.90,1.255) (12.00,1.000)};
    #1
  \end{tikzpicture}%
}

% Exercice 7 : f(x) = -x^2 + 2x + 8 et g(x) = 3(x-3)^2 - 3 sur [-3 ; 5]
\newcommand{\repereFG}[1]{%
  \begin{tikzpicture}[x=0.9cm, y=0.4cm]
    \axes{-3}{5}{-7}{10}
    \foreach \x in {-3,...,-1,1,2,...,5} \node[below] at (\x,0) {\scriptsize $\x$};
    \foreach \y in {-6,-4,-2,2,4,...,10} \node[left] at (0,\y) {\scriptsize $\y$};
    \node[below left] at (0,0) {\scriptsize $0$};
    \begin{scope}
      \clip (-3,-7) rectangle (5,10);
      \draw[very thick, domain=-3:5, samples=80, smooth] plot (\x, {-\x*\x+2*\x+8});
      \draw[very thick, dashed, domain=0.8:5, samples=80, smooth] plot (\x, {3*(\x-3)*(\x-3)-3});
    \end{scope}
    \node at (-1.3,4.5) {$\mathcal{C}_f$};
    \node at (4.6,7) {$\mathcal{C}_g$};
    #1
  \end{tikzpicture}%
}

% Exercice 8 : tarifs de location (A : 5x ; B : 15 + 2x)
\newcommand{\repereVelo}[1]{%
  \begin{tikzpicture}[x=0.8cm, y=0.1cm]
    \draw[gray!20, very thin, xstep=0.5, ystep=2.5] (0,0) grid (10,50);
    \draw[gray!45, very thin, xstep=1, ystep=5] (0,0) grid (10,50);
    \draw[-{Latex[length=2mm]}] (0,0) -- (10.5,0) node[above left] {\scriptsize durée (h)};
    \draw[-{Latex[length=2mm]}] (0,0) -- (0,53) node[above] {prix (€)};
    \foreach \x in {1,...,10} \node[below] at (\x,0) {\scriptsize $\x$};
    \foreach \y in {5,10,...,50} \node[left] at (0,\y) {\scriptsize $\y$};
    \node[below left] at (0,0) {\scriptsize $0$};
    \draw[very thick] (0,0) -- (10,50) node[pos=0.85, left] {A};
    \draw[very thick, dashed] (0,15) -- (10,35) node[pos=0.95, below] {B};
    #1
  \end{tikzpicture}%
}

% Exercice 9 : triangle ABC (AB = 4, AC = 8) et rectangle AMNP
\newcommand{\triangleABC}{%
  \begin{tikzpicture}[scale=0.6]
    \fill[gray!25] (0,0) rectangle (5,1.5);
    \draw[thick] (0,0) node[below left] {$A$} -- (8,0) node[below right] {$C$} -- (0,4) node[above left] {$B$} -- cycle;
    \draw[thick] (0,1.5) node[left] {$M$} -- (5,1.5) node[above right] {$N$} -- (5,0) node[below] {$P$};
    \draw[{Latex[length=1.5mm]}-{Latex[length=1.5mm]}] (-1,0) -- node[left] {$x$} (-1,1.5);
  \end{tikzpicture}%
}
\newcommand{\repereRect}[1]{%
  \begin{tikzpicture}[x=1.7cm, y=0.42cm]
    \draw[gray!20, very thin, xstep=0.25, ystep=0.5] (0,0) grid (4,9);
    \draw[gray!45, very thin, xstep=0.5, ystep=1] (0,0) grid (4,9);
    \draw[-{Latex[length=2mm]}] (0,0) -- (4.2,0) node[right] {$x$};
    \draw[-{Latex[length=2mm]}] (0,0) -- (0,9.5) node[above] {aire (cm$^2$)};
    \foreach \x/\l in {0.5/0{,}5,1/1,1.5/1{,}5,2/2,2.5/2{,}5,3/3,3.5/3{,}5,4/4} \node[below] at (\x,0) {\scriptsize $\l$};
    \foreach \y in {1,...,9} \node[left] at (0,\y) {\scriptsize $\y$};
    \node[below left] at (0,0) {\scriptsize $0$};
    \draw[very thick, domain=0:4, samples=60, smooth] plot (\x, {\x*(8-2*\x)});
    #1
  \end{tikzpicture}%
}

% Exercice 10 : rectangle ABCD (AB = 6, AD = 14), x = 2
\newcommand{\rectABCD}{%
  \begin{tikzpicture}[scale=0.42]
    \fill[pattern=north east lines] (0,0) rectangle (2,2);
    \fill[pattern=north east lines] (2,2) rectangle (14,6);
    \draw[very thick] (0,0) node[below left] {$A$} rectangle (14,6);
    \node[above left] at (0,6) {$B$};
    \node[above right] at (14,6) {$C$};
    \node[below right] at (14,0) {$D$};
    \draw[thick] (0,2) node[left] {$M$} -- (2,2) node[above left] {$N$} -- (2,0) node[below] {$P$};
    \draw[thick] (2,2) -- (2,6) node[above] {$J$};
    \draw[thick] (2,2) -- (14,2) node[right] {$I$};
    \node[left] at (0,1) {$x$};
    \node[below] at (1,0) {$x$};
  \end{tikzpicture}%
}
\newcommand{\repereABCD}[1]{%
  \begin{tikzpicture}[x=1.25cm, y=0.05cm]
    \draw[gray!20, very thin, xstep=0.25, ystep=2] (0,0) grid (6,90);
    \draw[gray!45, very thin, xstep=0.5, ystep=10] (0,0) grid (6,90);
    \draw[-{Latex[length=2mm]}] (0,0) -- (6.3,0) node[right] {$x$};
    \draw[-{Latex[length=2mm]}] (0,0) -- (0,95);
    \foreach \x in {1,...,6} \node[below] at (\x,0) {\scriptsize $\x$};
    \foreach \y in {10,20,...,90} \node[left] at (0,\y) {\scriptsize $\y$};
    \node[below left] at (0,0) {\scriptsize $0$};
    \draw[very thick, domain=0:6, samples=60, smooth] plot (\x, {2*\x*\x-20*\x+84});
    \draw[very thick, domain=0:6, samples=60, smooth] plot (\x, {-2*\x*\x+20*\x});
    #1
  \end{tikzpicture}%
}

% Exercice 11 : bénéfice B(x) = -2x^2 + 26x - 60 sur [3 ; 10]
\newcommand{\repereBallon}[1]{%
  \begin{tikzpicture}[x=0.72cm, y=0.15cm]
    \draw[gray!20, very thin, xstep=0.5, ystep=1] (0,0) grid (11,30);
    \draw[gray!45, very thin, xstep=1, ystep=5] (0,0) grid (11,30);
    \draw[-{Latex[length=2mm]}] (0,0) -- (11.4,0) node[above left] {\scriptsize prix (€)};
    \draw[-{Latex[length=2mm]}] (0,0) -- (0,31.5) node[above right] {\scriptsize bénéfice (milliers d'€)};
    \foreach \x in {1,...,11} \node[below] at (\x,0) {\scriptsize $\x$};
    \foreach \y in {5,10,...,30} \node[left] at (0,\y) {\scriptsize $\y$};
    \node[below left] at (0,0) {\scriptsize $0$};
    \draw[very thick, domain=3:10, samples=60, smooth] plot (\x, {-2*\x*\x+26*\x-60});
    #1
  \end{tikzpicture}%
}

\begin{document}
\pagestyle{fancy}

\fancyhf{} % clear existing header/footer entries
\fancyhead[L]{ {\Niveau}}
\fancyhead[C]{ {\Chapitre}}
\fancyhead[R]{ {\Document}}
\fancyfoot[L]{ A. Fernandez }
\fancyfoot[R]{{\thepage\  / \pageref{LastPage}}}

% =====================================================
\begin{exercice}{1}
\begin{minipage}{0.47\linewidth}
On a représenté ci-contre une fonction $f$.
\begin{enumerate}[label=\arabic*.]
  \item Quel est l'ensemble de définition de $f$~?
  \item Déterminer $f(-5)$, $f(2)$ et $f(6)$.
  \item Déterminer les antécédents de $0$, $5$ et $-3$.
  \item Le nombre $-4$ a-t-il des antécédents~?
  \item Compléter le tableau de valeurs.
\end{enumerate}
\end{minipage}
\hfill
\begin{minipage}{0.5\linewidth}
\begin{center}
\repereUn{}
\end{center}
\end{minipage}

\medskip
\begin{center}
\renewcommand{\arraystretch}{1.5}
\begin{tabular}{|c|*{6}{C{1.5cm}|}}
  \hline
  $x$ & $-7$ & $-3$ & $-1$ & $7$ & & \\ \hline
  $f(x)$ & & & & & $3$ & $0$ \\ \hline
\end{tabular}
\end{center}
\end{exercice}

% =====================================================
\begin{exercice}{2}
\begin{enumerate}[label=\arabic*.]
  \item Soit $f$ la fonction définie sur $\mathbb{R}$ par $f(x) = 3x - 5$.
  \begin{tasks}(2)
    \task Calculer $f(2)$, $f(-1)$, $f(0)$ et $f\left(\dfrac{1}{3}\right)$.
    \task Trouver l'antécédent de $7$, puis celui de $0$.
  \end{tasks}
  \item Soit $g$ la fonction définie sur $\mathbb{R}$ par $g(x) = x^2 - 4$.
  \begin{tasks}(2)
    \task Calculer $g(-3)$, $g(0)$ et $g\left(\dfrac{1}{2}\right)$.
    \task Déterminer les antécédents de $0$ par $g$.
    \task Le nombre $-5$ a-t-il des antécédents~? Justifier.
  \end{tasks}
\end{enumerate}
\end{exercice}

% =====================================================
\begin{exercice}{3}
On étudie la fonction $f$ définie sur $\mathbb{R}$ par $f(x) = 0{,}5x^2 - 0{,}5x - 3$.
\begin{enumerate}[label=\arabic*.]
  \item Compléter le tableau de valeurs.

  \medskip
  \renewcommand{\arraystretch}{1.5}
  \begin{tabular}{|c|*{8}{C{1.2cm}|}}
    \hline
    $x$ & $-3$ & $-2$ & $-1$ & $0$ & $1$ & $2$ & $3$ & $4$ \\ \hline
    $f(x)$ & & & & & & & & \\ \hline
  \end{tabular}
\end{enumerate}

\begin{minipage}{0.53\linewidth}
\begin{enumerate}[label=\arabic*., start=2]
  \item Placer les points dans le repère ci-contre, puis tracer la courbe de $f$ à main levée.
  \item $A(1\,;\,-3)$ appartient-il à la courbe de $f$~?
  \item Même question pour $B(5\,;\,6)$.
\end{enumerate}
\end{minipage}
\hfill
\begin{minipage}{0.45\linewidth}
\begin{center}
\repereVide{}
\end{center}
\end{minipage}
\end{exercice}

% =====================================================
\begin{exercice}{4}
Soit $f$ la fonction définie sur $\mathbb{R}$ par $f(x) = 2x^2 - 3x + 1$, et $\mathcal{C}_f$ sa courbe représentative.
\begin{enumerate}[label=\arabic*.]
  \item Les points $A(2\,;\,3)$, $B(-1\,;\,5)$ et $C\left(\dfrac{1}{2}\,;\,0\right)$ appartiennent-ils à $\mathcal{C}_f$~?
  \item Déterminer l'ordonnée du point $D$ de $\mathcal{C}_f$ d'abscisse $-2$.
  \item Déterminer les coordonnées du point où $\mathcal{C}_f$ coupe l'axe des ordonnées.
\end{enumerate}
\end{exercice}

% =====================================================
\begin{exercice}{5}
Elias consulte son compte bancaire. Le graphique représente le solde de son compte au cours de l'année.

Ici, $x$ est le nombre de mois écoulés depuis le 1\textsuperscript{er} janvier.

\begin{center}
\repereBanque{}
\end{center}

\begin{enumerate}[label=\arabic*.]
  \item Quel était le solde du compte le 1\textsuperscript{er} janvier~? Au bout de 10 mois~?
  \item À quels moments le solde était-il de $100$~€~?
  \item Pendant quelle période Elias était-il \og à découvert \fg{} (solde strictement négatif)~?
  \item Pendant quelle(s) période(s) le solde était-il supérieur ou égal à $100$~€~?
  \item Quel a été le solde le plus élevé de l'année~? À quel moment~?
\end{enumerate}
\end{exercice}

% =====================================================
\begin{exercice}{6}
On reprend la fonction $f$ de l'exercice 1.
\begin{enumerate}[label=\arabic*.]
  \item Résoudre graphiquement les équations suivantes.
  \begin{tasks}(4)
    \task $f(x) = 0$
    \task $f(x) = 3$
    \task $f(x) = 5$
    \task $f(x) = -4$
  \end{tasks}
  \item Résoudre graphiquement les inéquations suivantes (réponses sous forme d'intervalles).
  \begin{tasks}(4)
    \task $f(x) < 0$
    \task $f(x) \geqslant 0$
    \task $f(x) > 3$
    \task $f(x) \leqslant 5$
  \end{tasks}
  \item (Pour aller plus loin) Résoudre $f(x) \geqslant 3$. Que remarque-t-on~?
\end{enumerate}
\end{exercice}

% =====================================================
\begin{exercice}{7}
\begin{minipage}{0.45\linewidth}
Voici les courbes de deux fonctions $f$ et $g$.

Elles sont définies sur $[-3\,;\,5]$.

\medskip
Résoudre graphiquement :
\begin{enumerate}[label=\arabic*.]
  \item $f(x) = g(x)$
  \item $f(x) > 8$
  \item $g(x) \leqslant 0$
  \item $f(x) > g(x)$
  \item $f(x) \geqslant g(x)$
  \item $f(x) < g(x)$
  \item $f(x) \leqslant g(x)$
\end{enumerate}
\end{minipage}
\hfill
\begin{minipage}{0.52\linewidth}
\begin{center}
\repereFG{}
\end{center}
\end{minipage}
\end{exercice}

% =====================================================
\begin{exercice}{8}
On reprend la location de vélos de la fiche calcul, pour une durée de $x$ heures ($0 \leqslant x \leqslant 10$).

Le prix est de $5x$ euros avec le tarif A, et de $15 + 2x$ euros avec le tarif B.

\begin{center}
\repereVelo{}
\end{center}

Répondre aux questions \textbf{à l'aide du graphique}.
\begin{tasks}[label=\arabic*.](2)
  \task Prix de 4 heures avec chaque tarif~?
  \task Durée pour laquelle les tarifs sont égaux~?
  \task Durées où B est strictement moins cher~?
  \task Avec $30$~€, durée maximale avec chaque tarif~?
\end{tasks}
\end{exercice}

% =====================================================
\begin{exercice}{9}
\begin{minipage}{0.6\linewidth}
$ABC$ est rectangle en $A$, avec $AB = 4$~cm et $AC = 8$~cm.

\smallskip
On place $M$ sur $[AB]$, $N$ sur $[BC]$ et $P$ sur $[AC]$.

\smallskip
Le quadrilatère $AMNP$ est un rectangle.

\smallskip
On note $x = AM$ et $f(x)$ l'aire du rectangle $AMNP$, en cm$^2$.
\end{minipage}
\hfill
\begin{minipage}{0.36\linewidth}
\begin{center}
\triangleABC
\end{center}
\end{minipage}

\begin{enumerate}[label=\arabic*.]
  \item Quel est l'ensemble de définition de $f$~?
  \item On admet que $MN = 8 - 2x$. Montrer que $f(x) = x(8 - 2x)$.
  \item Calculer $f(1)$ et $f(1{,}5)$.
\end{enumerate}

\begin{enumerate}[label=\arabic*., start=4]
  \item On a tracé ci-dessous la courbe de $f$. Répondre à l'aide du graphique.

  \begin{center}
  \repereRect{}
  \end{center}

  \begin{enumerate}[label=\alph*.]
    \item Pour quelles valeurs de $x$ l'aire vaut-elle $6$~cm$^2$~?
    \item Pour quelles valeurs de $x$ l'aire est-elle d'au moins $6$~cm$^2$~?
    \item Où placer $M$ pour que l'aire soit la plus grande possible~? Que vaut-elle alors~?
  \end{enumerate}
  \item (Pour aller plus loin) Démontrer que $MN = 8 - 2x$ à l'aide du théorème de Thalès.
\end{enumerate}
\end{exercice}

% =====================================================
\newpage
\begin{exercice}{10}
\begin{minipage}{0.57\linewidth}
$ABCD$ est un rectangle tel que $AB = 6$ et $AD = 14$.

$M$ est un point du segment $[AB]$.

On construit le carré $AMNP$ et le rectangle $NICJ$.

On pose $AM = x$ et on note :
\begin{itemize}
  \item $f(x)$ l'aire de la partie hachurée ;
  \item $g(x)$ l'aire de la partie non hachurée.
\end{itemize}
\end{minipage}
\hfill
\begin{minipage}{0.4\linewidth}
\begin{center}
\rectABCD
\end{center}
\end{minipage}

\begin{enumerate}[label=\arabic*.]
  \item Dans quel intervalle varie $x$~?
  \item \begin{enumerate}[label=\alph*.]
    \item Exprimer $NI$ et $NJ$ en fonction de $x$, puis montrer que $f(x) = 2x^2 - 20x + 84$.
    \item En déduire que $g(x) = -2x^2 + 20x$.
  \end{enumerate}
\end{enumerate}

\begin{enumerate}[label=\arabic*., start=3]
  \item On a tracé ci-dessous les courbes de $f$ et de $g$. Laquelle représente $f$~? Justifier.

  \begin{center}
  \repereABCD{}
  \end{center}
  \item Où placer $M$ pour que l'aire hachurée soit égale à l'aire non hachurée~? Vérifier par le calcul.
  \item Pour quelle valeur de $x$ l'aire hachurée semble-t-elle la plus petite possible~?
\end{enumerate}
\end{exercice}

% =====================================================
\newpage
\begin{exercice}{11}
Une entreprise veut vendre des ballons de football. On note $x$ le prix d'un ballon, en euros.

Plus le prix est élevé, moins elle en vend. On cherche le prix qui donne le plus grand bénéfice.

\medskip
\textbf{Partie A.} Selon une étude, l'entreprise vend \mbox{$f(x) = 20 - 2x$} milliers de ballons, avec $x \in [0\,;\,10]$.
\begin{enumerate}[label=\arabic*.]
  \item Combien de ballons sont vendus si le prix est de $5$~€~? De $9$~€~?
  \item Quel prix faut-il choisir pour vendre $6$ milliers de ballons~? $12$ milliers~?
\end{enumerate}

\textbf{Partie B.} Fabriquer un ballon coûte $3$~€. On note $B(x)$ le bénéfice, en milliers d'euros.
\begin{enumerate}[label=\arabic*.]
  \item \begin{enumerate}[label=\alph*.]
    \item Expliquer pourquoi $B(x) = (20 - 2x)(x - 3)$.
    \item Montrer que $B(x) = -2x^2 + 26x - 60$.
  \end{enumerate}
  \item \begin{enumerate}[label=\alph*.]
    \item Calculer le bénéfice pour un prix de $6$~€.
    \item Que se passe-t-il si le prix est de $2$~€~?
  \end{enumerate}
  \item Compléter le tableau de valeurs par le calcul.

  \medskip
  \renewcommand{\arraystretch}{1.5}
  \begin{tabular}{|c|*{9}{C{1cm}|}}
    \hline
    $x$ & $1$ & $2$ & $3$ & $4$ & $5$ & $6$ & $7$ & $8$ & $9$ \\ \hline
    $B(x)$ & & & & & & & & & \\ \hline
  \end{tabular}
\end{enumerate}

\begin{enumerate}[label=\arabic*., start=4]
  \item On a représenté $B$ sur l'intervalle $[3\,;\,10]$. Répondre à l'aide du graphique.

  \begin{center}
  \repereBallon{}
  \end{center}

  \begin{enumerate}[label=\alph*.]
    \item Pour quels prix le bénéfice est-il supérieur ou égal à $20$ milliers d'euros~?
    \item Pour quels prix le bénéfice est-il strictement inférieur à $12$ milliers d'euros~?
    \item Quel prix donne le plus grand bénéfice~? Quel est ce bénéfice~?
  \end{enumerate}
\end{enumerate}
\end{exercice}
\end{document}
